\section*{第 \ref{ch:b2:weather-systems} 章——天气系统}

\begin{solution}{exo:b2:weather-systems:1}
(a) $f = 2 \times 7.292 \times 10^{-5} \times \sin\ang{50} = \qty{1.117e-4}{s^{-1}}$。(b) 取
$\bar T = \qty{278}{K}$，$\tan\alpha = 1.117 \times 10^{-4} \times 278 \times 15/(9.81 \times 6)
= 0.00791$，约为 1/126。(c) $1/0.00791 = 126$：锋后约 \qty{126}{km}。(d) 坡度
与温度差成反比，减半为 1/253。
\end{solution}

\begin{solution}{exo:b2:weather-systems:2}
(a) $6 \times 150 = \qty{900}{km}$，在锋前。(b) $900/25 = \qty{36}{h}$，即一天半。(c) 先
出现高空卷云，随后是卷层云，再是高层云；云层增厚、降低，
随\omterm{def:b2:weather-systems:fronts}{锋面}临近，自低而厚的云层开始降雨。
\end{solution}

\begin{solution}{exo:b2:weather-systems:3}
符号指示\omterm{def:b2:weather-systems:fronts}{锋面}移动的方向，并指向冷锋前方较暖的空气，或
暖锋前方较冷的空气。(a) 冷锋：冷空气位于其西北侧、锋后，
\omterm{def:b2:weather-systems:fronts}{锋面}向东南移动。(b) 暖锋：冷空气位于其北侧、锋前，\omterm{def:b2:weather-systems:fronts}{锋面}向
北移动。(c) \omterm{def:b2:weather-systems:fronts}{锢囚锋}连接低压与冷、暖锋起始处的三相点，
两侧均为冷空气，暖空气抬升至上方；\omterm{def:b2:weather-systems:fronts}{锋面}沿其符号所指的
方向移动。(d) \omterm{def:b2:weather-systems:fronts}{准静止锋}：三角形朝南指向暖空气，因此冷空气
位于北侧；\omterm{def:b2:weather-systems:fronts}{锋面}几乎不动。
\end{solution}

\begin{solution}{exo:b2:weather-systems:4}
(a) $\beta = 2 \times 7.292 \times 10^{-5} \times \cos\ang{45}/(6.371 \times 10^6) =
\qty{1.62e-11}{m^{-1}.s^{-1}}$。(b) 对于 \qty{4000}{km}，$k = \qty{1.571e-6}{m^{-1}}$，$\beta/k^2 =
\qty{6.6}{m/s}$，且 $c = 12 - 6.6 = \qty{5.4}{m/s}$；对于 \qty{8000}{km}，$\beta/k^2 =
\qty{26.3}{m/s}$，且 $c = \qty{-14.3}{m/s}$。(c) 短波向东移动，每天 \qty{470}{km}；
长波向西移动，每天 \qty{1240}{km}，逆风而行。
\end{solution}

\begin{solution}{exo:b2:weather-systems:5}
(a) 当 $c = 0$ 时，$U = \beta/k^2$，故 $k = \sqrt{\beta/U}$，且 $\lambda = 2\pi/k =
2\pi\sqrt{U/\beta}$。(b) $2\pi\sqrt{12/(1.62 \times 10^{-11})} = \qty{5410}{km}$。(c) 在纬度
\ang{45} 处，纬圈长 $2\pi \times 6371 \times \cos\ang{45} = \qty{28300}{km}$，约 5.2
个波长，实际取 5 个。(d) $\lambda$ 随 $\sqrt U$ 增大：较强的西风使
驻波变长，绕半球的波数减少。
\end{solution}

\begin{solution}{exo:b2:weather-systems:6}
(a) 在 $C_+ = (\sqrt{72}, \sqrt{72}, 27)$ 处：$\dot x = 10(\sqrt{72} - \sqrt{72}) = 0$，$\dot y =
\sqrt{72}(28 - 27) - \sqrt{72} = 0$, $\dot z = 72 - \tfrac83 \times 27 = 72 - 72 = 0$.
(b) $\mathrm{e}^{-13.67} = 1.16 \times 10^{-6}$：体积在一个
时间单位内缩小近百万倍。(c) 体积可以缩小，即在某一方向上被拉伸，而
在其他方向上的压缩要强得多：沿拉伸方向的距离增大，两条邻近轨迹
沿此方向分离，而它们之间的体积趋于零。
\end{solution}

\begin{solution}{exo:b2:weather-systems:7}
(a) $\ln 2/0.906 = 0.765$ 个时间单位。(b) $\ln(10/10^{-8})/0.906 = \ln(10^9)/0.906 = 22.9$ 个
时间单位。(c) $\ln 1000/0.906 = 7.6$ 个时间单位：初始状态改善千倍，
只换得 (b) 中时间的三分之一，因为 $\ln 10^3/\ln 10^9 = 1/3$。
\end{solution}

\begin{solution}{exo:b2:weather-systems:8}
(a) $(303 - 195)/195 = 0.554$，且 $V_{\max} = \sqrt{0.9 \times 0.554 \times 15\,000} =
\qty{86.5}{m/s}$。(b) 远高于 5 级的 \qty{70}{m/s}，接近
近期五场风暴曾超过的 \qty{86}{m/s}。(c) 风暴很少长时间处于理论所假设的
稳态；垂直风切变、被卷入风暴的干燥空气以及周围的下沉气流，
这些理论未予考虑的因素会削弱风暴。
\end{solution}

\begin{solution}{exo:b2:weather-systems:9}
(a) $(303 - 195)/303 = 0.356$。(b) $(303 - 195)/195 = 0.554$；比值为 $0.554/0.356 =
303/195 = 1.554$。(c) 在卡诺热机中，功离开热机。在热带气旋中，
功因海面附近的摩擦而耗散，这部分热量加热了温度
为 $T_s$ 的流入空气，即热机吸热之处；热量被再次利用，热机表现
如同效率提高了 $T_s/T_o$ 倍。
\end{solution}

\begin{solution}{exo:b2:weather-systems:10}
(a) 不加近似地由\cref{prop:b2:weather-systems:slope}推导，可得
$\mathrm{d}z/\mathrm{d}x = f(\rho_c v_c - \rho_w v_w)/[g(\rho_c - \rho_w)]$。密度较大的流体
静置于较轻流体之上是不稳定的：哪怕最微小的扰动也会使其翻转。在
静止的\omterm{def:b2:weather-systems:fronts}{锋面}中，冷而密度较大的空气必定位于下方，这正是定义中的楔形；于是
$\rho_c > \rho_w$，界面向冷空气一侧抬升，且 $\rho_c v_c > \rho_w v_w$。
(b) $\zeta = \partial v/\partial x = 10/(50 \times 10^3) = \qty{2e-4}{s^{-1}}$，约为 $1.9f$，其中
$f = \qty{1.03e-4}{s^{-1}}$。(c) 北半球为正，呈气旋式，大小与
科里奥利参数本身相当；\omterm{def:b2:weather-systems:fronts}{锋面}是强气旋式涡度的带。
\end{solution}

\begin{solution}{exo:b2:weather-systems:11}
(a) $\log_2(6/0.5) = \log_2 12 = 3.58$ 次倍增，$3.58 \times 2.4 = 8.6$~天。(b) 若
多预报一天，则误差在失效前需少增长一个倍增时间，初始误差
须每个十进位缩小 $2^{1/2.4} = 1.33$ 倍。(c) 缩小十倍相当于
$\log_2 10 = 3.32$ 次倍增；若仅增加 4 至 5 天，则小误差倍增一次所需时间为 $4/3.32 = 1.2$
至 $5/3.32 = 1.5$ 天：小误差的增长快于预报中 20 世纪 80 年代所测的较大误差，
这是由小尺度对流释放潜热所致。
\end{solution}

\begin{solution}{exo:b2:weather-systems:12}
(a) 在 \qty{30}{\celsius} 时，$e_s = \qty{43.5}{hPa}$，$q_s^* = \qty{27.1}{g/kg}$；空气在
\qty{29}{\celsius} 时，$e = 0.8 \times 41.0 = \qty{32.8}{hPa}$，$q_a = \qty{20.4}{g/kg}$；
$h_o^* - h^* = 1004.5 \times 1 + 2.501 \times 10^6 \times 0.0067 = \qty{17.8}{kJ/kg}$，对照
\qty{16.8}{kJ/kg}。(b) $(303.15 - 200)/200 = 0.516$，且 $V_{\max} = \sqrt{0.9 \times 0.516
\times 17\,800} = \qty{90.9}{m/s}$，对照 \qty{87.9}{m/s}。(c) 约每
开尔文 \qty{3.4}{\percent}。非平衡度升高 \qty{6}{\percent}，使速度增加 \qty{3}{\percent}；
因子 $(T_s - T_o)/T_o$ 升高 \qty{1}{\percent}，使速度增加 \qty{0.5}{\percent}：贡献最大的
是燃料，即较暖海面所能提供的水汽，它通过 $e_s$ 的
指数式增长起作用。
\end{solution}

\begin{solution}{pb:b2:weather-systems:1}
\textbf{1.} $(480 - 390)/6 = \qty{15}{km/h}$.
\textbf{2.} $390/15 = \qty{26}{h}$：在 12 UTC 之后，即次日 14 UTC。
\textbf{3.} $f = \qty{1.117e-4}{s^{-1}}$.
\textbf{4.} 取 $\bar T = \qty{281}{K}$，$\tan\alpha = 1.117 \times 10^{-4} \times 281 \times
16/(9.81 \times 8) = 0.0064$，约为 1/156。
\textbf{5.} 此时\omterm{def:b2:weather-systems:fronts}{锋面}位于镇外 $6 \times 15 = \qty{90}{km}$ 处，该处地表比镇子
高出 $90\,000 \times 0.0064 = \qty{576}{m}$。
\textbf{6.} 位于锋后 $1500/0.0064 = \qty{234}{km}$ 处，过境后
$234/15 = \qty{15.6}{h}$ 到达。
\textbf{7.} 积云与积雨云、阵雨，或伴有雷声、阵风，以及气温
骤降；气压在\omterm{def:b2:weather-systems:fronts}{锋面}临近时下降，过境时开始回升；随后
为较凉爽、较清澈的空气。
\textbf{8.} $\beta = 2 \times 7.292 \times 10^{-5} \times \cos\ang{50}/(6.371 \times 10^6) =
\qty{1.47e-11}{m^{-1}.s^{-1}}$.
\textbf{9.} $k = 2\pi/(5 \times 10^6) = \qty{1.257e-6}{m^{-1}}$，且 $\beta/k^2 = \qty{9.3}{m/s}$。
\textbf{10.} $c = 20 - 9.3 = \qty{10.7}{m/s}$，向东。
\textbf{11.} $10.7 \times 86\,400 = \qty{924}{km}$，向东。
\textbf{12.} $\lambda_s = 2\pi\sqrt{20/(1.47 \times 10^{-11})} = \qty{7330}{km}$；纬圈
在 \ang{50} 处的纬圈长 $2\pi \times 6371 \times \cos\ang{50} = \qty{25700}{km}$，约
相当于 3.5 个这样的波。
\textbf{13.} 在其前方（即东侧），短波在此诱发上升气流，
地面气压下降。
\textbf{14.} $\lambda = \ln(3.9 \times 10^{-3}/5 \times 10^{-7})/(14 - 4) = \ln 7800/10 = 0.896$.
\textbf{15.} $\ln 2/0.896 = 0.77$ 个时间单位。
\textbf{16.} 比 0.906 低 \qty{1}{\percent}：200 对的平均值、有限时间及
从图中读取的数值。
\textbf{17.} 两次运行一旦相距甚远，它们便位于吸引子上互不相关的点，且
它们的距离不能超过其大小：增长趋于饱和。
\textbf{18.} $\ln 10/0.896 = 2.6$ 个时间单位之后。
\textbf{19.} $\log_2(3/0.1) = \log_2 30 = 4.91$ 次倍增。
\textbf{20.} $4.91 \times 2.4 = 11.8$~天。
\textbf{21.} 一个倍增时间：2.4~天。
\textbf{22.} $4.91 \times 5 = 24.5$~天：早先的估计过于乐观。
\textbf{23.} 倍增时间为 2.4 天，若误差缩小十倍，则可增加 $3.32 \times 2.4
= 8.0$~天，而非 4 至 5 天：小误差每次倍增所需时间少于 2.4 天，
据此作出的估计高估了预报时效。
\textbf{24.} 该预报至多约可用 \textbf{12 天}，实际上
时效还要更短。
\end{solution}
